\chapter{Bistochastic Matrices and the Birkhoff Polytope}
\label{ch:birkhoff}

\section{Bistochasticity}

A stochastic map $M$ (Chapter~\ref{ch:simplex}) is \emph{bistochastic} (doubly
stochastic) if, in addition to column sums, its row sums also equal one:
\begin{equation}
B_{ij} \ge 0, \qquad \sum_i B_{ij} = 1, \qquad \sum_j B_{ij} = 1.
\end{equation}
Write $\mathrm{DS}(n)$ for the set of $n \times n$ bistochastic matrices. Row
normalization is the additional admissibility constraint motivating this
chapter: a bistochastic map sends the uniform distribution to itself, and (as
Theorem~\ref{thm:birkhoff-vonneumann} below shows) is built entirely from
maps that merely relabel alternatives.

\section{The Birkhoff Polytope}

Bistochastic matrices have long been treated as the natural probabilistic
analogue of reversible transformations: permutations are the only
transformations on a finite set that never lose distinguishability, and a
bistochastic matrix is, by the theorem below, always assembled from a mixture
of them. That association motivates why this particular restriction of
$\mathrm{DS}(n)$ — not some other subset — is the one worth naming and
studying before narrowing the admissible class further.

\begin{definition}
A \emph{permutation matrix} $P_\sigma$, for $\sigma \in S_n$, has
$(P_\sigma)_{ij} = \delta_{i,\sigma(j)}$. Every permutation matrix is
bistochastic.
\end{definition}

\begin{theorem}[Birkhoff--von Neumann]
\label{thm:birkhoff-vonneumann}
\index{Birkhoff--von Neumann (theorem)}
$\mathrm{DS}(n)$ is a convex polytope (the \emph{Birkhoff polytope}) whose
extreme points are exactly the $n!$ permutation matrices. Consequently every
$B \in \mathrm{DS}(n)$ can be written as a convex combination
\begin{equation}
B = \sum_{\sigma \in S_n} c_\sigma P_\sigma, \qquad c_\sigma \ge 0, \qquad
\sum_\sigma c_\sigma = 1,
\end{equation}
though not, in general, uniquely.
\end{theorem}

\begin{proof}
$\mathrm{DS}(n)$ is defined by finitely many linear equalities (row/column
sums) and inequalities ($B_{ij}\ge0$) in a bounded region of $\R^{n^2}$, hence
is a convex polytope.

\emph{Permutation matrices are extreme.} Suppose $P_\sigma = \tfrac12(M_1+M_2)$
for $M_1,M_2 \in \mathrm{DS}(n)$. Wherever $(P_\sigma)_{ij}=0$, nonnegativity of
$M_1,M_2$ forces $(M_1)_{ij}=(M_2)_{ij}=0$. Wherever $(P_\sigma)_{ij}=1$, each
row of $M_1$ (resp.\ $M_2$) sums to $1$ with all entries $\ge0$, so
$(M_1)_{ij}\le1$ and likewise for $M_2$; averaging to $1$ with both terms
$\le1$ forces $(M_1)_{ij}=(M_2)_{ij}=1$. Hence $M_1=M_2=P_\sigma$.

\emph{Every extreme point is a permutation matrix.} Let $B \in \mathrm{DS}(n)$
be arbitrary and not a permutation matrix, and let $G_B$ be the bipartite
support graph on row-vertices and column-vertices with an edge $(i,j)$
whenever $B_{ij}>0$. Every vertex of $G_B$ has degree $\ge1$ (each row/column
sums to $1>0$), and since $B$ is not a permutation matrix, some vertex has
degree $\ge2$. A finite graph with a vertex of degree $\ge2$ and no vertex of
degree $0$ contains a cycle; since $G_B$ is bipartite, this cycle has even
length and alternates rows and columns,
\begin{equation}
i_1,\ j_1,\ i_2,\ j_2,\ \dots,\ i_k,\ j_k,\ i_1, \qquad
B_{i_l j_l} > 0, \quad B_{i_{l+1} j_l} > 0.
\end{equation}
Let $\varepsilon > 0$ be smaller than every entry of $B$ along the cycle, and
define $\Delta$ by alternating $+\varepsilon,-\varepsilon$ along the cycle
edges (zero elsewhere). Because the pattern alternates sign around a closed
cycle, every row and column sum of $\Delta$ is zero, so $B \pm \Delta$ are
again nonnegative (by choice of $\varepsilon$) with the same row and column
sums as $B$: both lie in $\mathrm{DS}(n)$. Since $B = \tfrac12\big[(B+\Delta) +
(B-\Delta)\big]$ and $\Delta \neq 0$, $B$ is not extreme. Contrapositively,
every extreme point is a permutation matrix.

The final claim (every $B \in \mathrm{DS}(n)$ is a convex combination of
extreme points) is the Krein--Milman property of a finite-dimensional compact
convex set applied to the polytope $\mathrm{DS}(n)$.
\end{proof}

\begin{remark}
Theorem~\ref{thm:birkhoff-vonneumann} already illustrates the book's recurring
pattern: a physically motivated constraint (bistochasticity) carves out a
polytope whose extreme structure (permutations) is exactly the reversible,
information-preserving special case, with everything else in the polytope
representable as a mixture of reversible maps. This is the classical precursor
to the much sharper restriction imposed by orthostochasticity and
unistochasticity in Chapter~\ref{ch:unistochastic}, where the analogous extreme
structure is no longer merely a finite set of permutations but a continuum
requiring genuine geometric (and, past $n=2$, non-convex) analysis.
\end{remark}

\section{Majorization}

Bistochastic maps have a further characterization in terms of \emph{mixing}:
applying one can only make a distribution more uniform, never less.

\begin{definition}[Majorization]
For $p, q \in \Delta_{n-1}$, write $q \prec p$ (``$q$ is majorized by $p$'')
if, after sorting both in decreasing order,
\begin{equation}
\sum_{i=1}^k q_i^{\downarrow} \le \sum_{i=1}^k p_i^{\downarrow} \qquad \text{for
every } k = 1,\dots,n-1
\end{equation}
(with equality automatic at $k=n$).
\end{definition}

\begin{theorem}[Hardy--Littlewood--P\'olya]
\label{thm:hlp}
\index{Hardy--Littlewood--P\'olya (theorem)}
$q \prec p$ if and only if $q = Bp$ for some $B \in \mathrm{DS}(n)$.
\end{theorem}

\begin{proof}
($\Leftarrow$) Suppose $q = Bp$ for $B \in \mathrm{DS}(n)$. Fix $k \in
\{1,\dots,n\}$ and let $S$ be \emph{any} set of $k$ indices. Then
\begin{equation}
\sum_{i \in S} q_i = \sum_{i \in S} \sum_j B_{ij} p_j = \sum_j w_j p_j,
\qquad w_j := \sum_{i \in S} B_{ij} \in [0,1], \qquad \sum_j w_j = k,
\end{equation}
the last equality because each column of $B$ sums to $1$ and $|S|=k$. Among
all weight vectors $w$ with entries in $[0,1]$ summing to $k$, $\sum_j w_j p_j$
is maximized by placing weight $1$ on the $k$ largest entries of $p$ (an
elementary rearrangement argument), giving
\begin{equation}
\sum_{i \in S} q_i \le \sum_{i=1}^k p_i^\downarrow
\end{equation}
for every $S$ of size $k$; in particular for $S$ the indices of the $k$
largest entries of $q$, this is exactly $\sum_{i=1}^k q_i^\downarrow \le
\sum_{i=1}^k p_i^\downarrow$. Since $k$ was arbitrary, $q \prec p$.

($\Rightarrow$) Suppose $q \prec p$. A \emph{$T$-transform} on indices
$(i,j)$ with parameter $t \in [0,1]$ is the bistochastic matrix
$T = I - t(e_i-e_j)(e_i-e_j)^\top$, a convex combination $t\,P_{(ij)} +
(1-t)\,I$ of the identity and the transposition $(i\,j)$; applying it moves
mass between coordinates $i,j$ while preserving their sum, replacing $(p_i,
p_j)$ by a pair strictly between them. The classical construction (Muirhead)
produces $q$ from $p$ by a finite sequence of such transfers, each moving
mass from a coordinate where $p$ (or its current partially-transformed image)
exceeds the corresponding partial sums of $q$ to one where it falls short,
terminating once all partial-sum inequalities defining $q \prec p$ are
saturated. Each $T$-transform lies in $\mathrm{DS}(n)$, and $\mathrm{DS}(n)$
is closed under matrix multiplication (row/column sums of $1$ and
nonnegativity are both preserved by composing two such maps), so the composite
$B = T_m \cdots T_1$ of the finitely many transforms used again lies in
$\mathrm{DS}(n)$, and $q = Bp$ by construction.
\end{proof}

\begin{corollary}[Entropy is not decreased by mixing]
\label{cor:entropy-increase}
\index{Entropy is not decreased by mixing (corollary)}
If $q = Bp$ for $B \in \mathrm{DS}(n)$, then $H(q) \ge H(p)$, where $H(p) =
-\sum_i p_i \log p_i$ is the Shannon entropy.
\end{corollary}

\begin{remark}
Corollary~\ref{cor:entropy-increase} follows from Theorem~\ref{thm:hlp} together
with the Schur-concavity of $H$: entropy increase under bistochastic evolution
is a \emph{structural} consequence of majorization geometry, not a
thermodynamic postulate added on top of the probabilistic framework. This is
the classical, pre-quantum ancestor of the entropy behavior discussed later
under coarse-graining and decoherence.
\end{remark}

\begin{ledgerentry}[End of Chapter~\ref{ch:birkhoff}]
\textbf{Established:} bistochastic matrices form a polytope generated by
permutations (Theorem~\ref{thm:birkhoff-vonneumann}); bistochastic evolution is
exactly majorization-compatible evolution (Theorem~\ref{thm:hlp}), and cannot
decrease Shannon entropy (Corollary~\ref{cor:entropy-increase}).\\
\textbf{Not yet established:} any restriction singling out orthostochastic or
unistochastic matrices within $\mathrm{DS}(n)$ — that restriction, and its
non-convex geometry for $n \ge 3$, is the subject of
Chapter~\ref{ch:unistochastic}.
\end{ledgerentry}
